% commentaire.tex — f176v, vv. 1169–1210 (conversion initiale des travaux antérieurs)
\comm{1170}{τὰ μὲν παρ᾽ ἡμῶν ἴσθι σοι πεπεισμένα·}{Une scholie intérieure au texte nous précise :\citComm{Ἀττικῶς ἀντὶ τοῦ `πεπείσμεθα'.}}
\comm{1174}{πρῶτον μὲν οὖν δίελθε κἀνακόλπασον.}{Pour le texte et le sens d'ἀνακαλπάζειν, voir le commentaire critique.}
\comm{1175}{Τερηδὼν,}{ἡ αὐλητρίς. (Note du scholiaste ; pourquoi pense-t-il que Térédon soit \italique{une} joueuse de flûte ?)}
\comm{1175}{Περσικόν.}{Euripide demande au flûtiste d'accompagner la danseuse avec un air \italique{persan} -- les oreilles ? Or il existe une danse du même nom\nc{Hésychius (π 2005) : περσικά· [ ... ] καὶ ὄρχησις Περσική.} que Xénophon nous décrit dans son \italique{Anabase} (VI, 1, 10) : \citComm{μετὰ τοῦτο Μυσὸς\nc{Je cite longuement le passage de Xénophon pour faire apparaître la race du danseur. Dans le cas où la danse persique aurait effectivement un rapport avec le verbe ἀνακαλπάζειν du vers 1174 -- ce que je ne crois pas --, on notera qu'il faisait partie du peuple d'Asie qui servit de titre à la pièce d'Eschyle, citée par Photius à l'adresse ἀνακαλπάζει et dont on a parlé au vers précédent.} εἰσῆλθεν ἐν ἑκατέρᾳ τῇ χειρὶ ἔχων πέλτην, καὶ τοτὲ μὲν ὡς δύο ἀντιταττομένων μιμούμενος ὠρχεῖτο, τοτὲ δὲ ὡς πρὸς ἕνα ἐχρῆτο ταῖς πέλταις, τοτὲ δ᾽ ἐδινεῖτο καὶ ἐξεκυβίστα ἔχων τὰς πέλτας, ὥστε ὄψιν καλὴν φαίνεσθαι. τέλος δὲ \textbf{τὸ περσικὸν ὠρχεῖτο} κρούων τὰς πέλτας καὶ ὤκλαζε καὶ ἐξανίστατο· καὶ ταῦτα πάντα ἐν ῥυθμῷ ἐποίει πρὸς τὸν αὐλόν.\\
Un Mysien entra ensuite ; il tenait, dans chacune de ses mains, un bouclier léger ; quelquefois il s'en servait en dansant, comme s'il eût eu à se défendre contre deux adversaires ; quelquefois, comme s'il n'eût eu affaire qu'à un seul. Souvent il tournait et faisait le saut périlleux sans lâcher ses boucliers : c'était un spectacle agréable qu'il donnait ; il finit par \textbf{danser la Persique} ; frappant d'un bouclier sur l'autre ; il s'accroupissait, se relevait, et exécutait tous ces mouvements en mesure et au son de la flûte.}L'accroupissement (ὀκλάζω) est le mouvement principal de cette danse et lui a donné son nom ; c'est le scholiaste aux \italique{Thesmophories} qui nous l'apprend : \citComm{Περσικὸν τὸ ὄκλασμα καλεῖται.}Pollux (IV, 100), compulsant un peu vite les \italique{Thesmophories} et ses commentateurs, croit que le terme ὄκλασμα apparaît dans le texte-même d'Aristophane : il range l'ὄκλασμα auprès du κνισμός parmi les danses guerrières et frénétiques :\citComm{καὶ ὄκλασμα· οὕτω γὰρ ἐν Θεσμοφοριαζούσαις ὀνομάζεται τὸ ὄρχημα τὸ Περσικὸν καὶ σύντονον. τὴνδ' αὐτὴν καὶ ὑγρὰν ὠνόμαζον.}L'on trouve d'ailleurs l'une des scholies des \italique{Thesmophories} les plus savantes à notre vers 1175\nc{Elle est discutée par Jacques Moreau dans l'article déjà mentionné : «~Note sur le vers 1174 des \italique{Thesmophories} d'Aristophane~», \italique{L'antiquité classique}, tome 7, fasc. 2 (1938) pp. 227-231.} :\citComm{Περσικὸν ὄρχημα, περὶ οὗ Ἰόβας μακρὸν πεποίηται λόγον ἐν τοῖς περὶ τῆς θεατρικῆς ἱστορίας. ὥστε λελύσθαι τὴν Σελεύκου πρότασιν. προτείνει γὰρ ἐν τῷ πρὸς Ζήνωνα προτατικῷ τὸ Περσικὸν ὄρχημα.}Pour ma part, je doute que l'exercice auquel se livre notre danseuse Ἐλάφιον soit cette danse guerrière illustrée par le Mysien de Xénophon ; l'adjectif Περσικός était assez commun pour qu'il pût désigner à la fois un air de flûte et une forme de danse sans qu'il y eût nécessairement de rapport entre les deux.}
\comm{1172 et 1180}{ὦλάφιον / ὡς ἐλαπρός, ὥσπερ ψύλλο κατὰ τὸ κῴδιο.}{Les verbes signifiant \italique{sauter}, \italique{danser} (πηδᾶν, ἄλλεσθαι) attirent presque inévitablement une comparaison, et le plus souvent avec un animal.
Les jeunes filles de Sparte dansent en bondissant comme des pouliches, ἇτε πῶλοι ταὶ κόραι / πὰρ τὸν Εὐρώταν / ἀμπαδίοντι (ἀμπαδίοντι Coulon : ἀμπάλλοντι codd.) (\italique{Lysistr.} 1308)\nc{Comparaison banale, cf. Eur. \italique{Or}. 45 (à propos d'Oreste furieux): πηδᾷ δρομαῖος, πῶλος ὣς ὑπὸ ζυγοῦ ; Eubul. 75, 6 $[=$ Athénée, III, 108ab$]$. Le πωλίον symbolise la vivacité, Ar. \italique{Guêp}. 188. En outre, πῶλος désigne couramment les jeunes femmes : Anacréon 88 \italique{Diehl} [$=$ 417 \italique{Poetae Melici Graeci} Page$]$ ; Eur. \italique{Héc.} 142, \italique{Hipp}. 516, etc. ; Cratin. 87 $[=$ 94 Kassel-Austin$]$ ; Épicrat. 9, 4; Eubul. 84 $[=$ 82 Kassel-Austin$]$.}. 
La Muse laconienne danse et saute comme une biche, ποδοῖν τε πάδη / ἇ τις ἔλαφος (\italique{ib}. 1317)\nc{Le cerf (ἔλαφος) est le symbole de l'agilité, \italique{Nuées} 351. (Et aussi, comme chez Homère, de la lâcheté.)}.
La danseuse des \italique{Thesmophories} porte d'ailleurs le sobriquet de \italique{Bichette} (Ἐλάφιον, v. 1172).
Bichette danse avec tant de grâce devant l'archer Scythe qu'il s'écrie : « qu'elle est légère ! une vraie puce sur ma peau de mouton » (\italique{Thesm}. 1180). Cette comparaison dénuée de poésie sent le corps de garde et ses peaux de mouton où grouille la vermine (on dormait communément sur des \italique{peaux de mouton}, κῴδια). (Note de Taillardat, §791.)}
\comm{1181}{φέρε θοἰμάτιον ἄνωθεν.}{Pour le texte et le sens (note 28 de Willems), voir le commentaire critique.}
\comm{1185}{ὡς στεριπὸ τὸ τιττί᾽}{$=$ ὡς στέριφον τὸ τιτθίον.}
\comm{1187}{κλαῦσί γ᾽ ἂν μὴ 'νδον μένῃς.}{πρὸς τὸ πέος λέγει ἐπιπλήττων ἐντεινομένῳ. (Note du scholiaste.)\\
L'archer suit d'un œil émérillonné tous les mouvements de la danseuse. Son agitation se trahit d'une manière par trop visible. Le Scythe, qui s'en aperçoit, reprend « gare à toi, etc. »}
\comm{1192}{ὠς γλυκερὸ τὸ γλῶσσ᾽, ὤσπερ Ἀττικὸς μέλις.}{Si l'image des baisers plus doux que le miel\nc{Et l'Attique est souvent citée pour son miel. Cf. la \italique{Paix}, vers 252-254 ; Antiphane fr. 177(179) Kassel-Austin $=$ Athénée II, 43b \& III, 74d ; le corpus hippocratique mentionne maintes fois du miel Attique.} ou le nectar nous paraît banale, c'est à une longue tradition héritée de ces culs-terreux de latins que nous le devons\nc{Cf. le \italique{Dictionnaire des sentences latines et grecques} de Renzo Tosi, §1385 \latin{Nectar dulcius} : Catulle 99, 1-2 ; Horace, \italique{Carmina} 1, 13, 14-16 ; \&c. Honoré de Balzac a fameusement repris cette longue tradition : « Les baisers d'une femme sincère ont un miel divin qui semble mettre dans cette caresse une âme, un feu subtil par lequel le cœur est pénétré. Les baisers dénués de cette onction savoureuse sont âpres et secs. »} ; pour ce qui est des auteurs de langue grecque, on ne peut guère citer qu'Alciphron (I, 38, 7), postérieur d'environ six cents ans à Aristophane. 
Les spectateurs d'Aristophane devaient ressentir avec cette comparaison une confusion burlesque : l'épithète poétique μελίγλωσσος existe en grec, mais l'image d'une \italique{langue pleine de miel} signifie tout autre chose. 
Ainsi parlait Taillardat (§§739-743) :\\
«~Le miel étant le symbole de la douceur\nc{\italique{Guêpes} 878 (voir §361), \italique{Thesm.} 1192, fr. 587, 622 $[=$ 681, 633 K.-A.$]$ ; Alexis 146, 6 $[=$ 150.6 K.-A.$]$ ; Ménand. 708 (?), etc.}, les poètes se plaisent à dire que leurs chants sont de miel\nc{Pind. \italique{Ol.} 11, 4 : μελιγάρυες ὕμνοι ; id. \italique{Ném.} 11,18 : μελιγδούποισι ... μελιζέμεν ἀοιδαῖς ; id. \italique{Isthm.} 2, 32 : μελικόμπων ἀοιδᾶν ; Alcman 41 $[=$ Athénée XIV, 633a $=$ fr. adesp. 954b \italique{Poetae Melici Graeci}, ed. Page$]$ : μέλεα μελιπτέρωτα ; Simonid. 40 D $[=$ fr 595.3 \italique{Poetae Melici Graeci} Page$]$ : μελιαδέα γᾶρυν, etc. On notera que, déjà chez Homère (\italique{Il.} I, 249), la voix de Nestor est dite $[$ἀπὸ γλώσσης$]$ μέλιτος γλυκίων. D'autre part Pindare, \italique{Ol.} 7, 8, appelle encore la poésie un nectar, νέκταρ.}. Cette image est chez Aristophane : μελιγλώσσων ἐπέων ... ἀοιδάν (\italique{Ois}. 908) donne un exemple de plus de l'épithète banale μελίγλωσσος\nc{Eschyl. \italique{Prom.} 172 ; Bacchyl. 3, 97 ; fr. 4, 25 Snell (μελιγλώσσων ἀοιδᾶν ἄνθεα).} $[$ ... $]$
On dira donc tout naturellement que la lèvre d'un poète est enduite de miel : c'est le cas de la bouche de Sophocle, Σοφοκλέους τοῦ μέλιτι κεχριμένου... τὸ στόμα (\italique{fr}. Aristophane 581 $[=$ 598 Kassel-Austin$]$). Composant des poèmes doux comme le miel, les poètes aiment à se comparer aux abeilles, ainsi qu'en témoigne Platon (\italique{Ion}, 534b).~»\\
Aristophane fait rire son public aux dépens de ces images trop usées : le Scythe rustique vante la \italique{langue douce comme le miel} de la danseuse, mais c'est une langue toute matérielle -- et qui appartient d'ailleurs à une actrice muette. Il y a des impropriétés de termes qui sont des malpropretés du cœur.}
\comm{1195}{ἐμοὶ κάρισο σὺ τοῦτο}{$=$ ἐμοὶ χάρισαι σὺ τοῦτο.}
\comm{1197–216}{ἀλλὰ τὸ συμβήνην λαβέ $[$ ... $]$ ὀρτῶς δὲ συβήνη 'στί· καταβηνῆσι γάρ.}{Ce passage a été si maltraité par les commentateurs -- anciens comme modernes -- que je me suis cru obligé à commettre cette note rectificatrice. \italique{Recte}, descendons-donc dans l'arène (ὀρτῶς δὲ συμβήνη 'στί· καταβηνῆσι γάρ.)\\
L'origine de tous les malentendus est la glose d'un commentateur ancien qui sera allé un peu vite : le \italique{Ravennas} comporte deux scholies pour définir la συμβήνην (la première scholie se trouve à l'intérieur même du texte) :\citComm{τὴν τοξοθήκην\\
συβήνη· αὐλοθήκη. λέγουσι δὲ καὶ τὸν φαρετρεῶνα συβήνην.}On retrouve cette scholie chez Hésychius (σ 2136) :\citComm{συβήνη· αὐλοθήκη, ἢ τοξοθήκη. ἢ ὁ ναυτικὸς χιτών.}Suidas cite notre passage dans une glose qu'on croise à l'identique, ou peu s'en faut, dans d'autres lexiques\nc{Photius (σ 660), \latin{Etymologicum magnum}, \latin{genuinum}, Zonaras, \&c. Pollux (VII, 153 ; X, 153) : καὶ τὸ μὲν ἀγγεῖον τὸ τὰς γλώττας ἔχον γλωττοκομεῖον, ἣ δὲ τῶν αὐλῶν θήκη συβήνη. Pour un dossier plus complet, je renvoie le lecteur à la \allemand{Realencyclopädie} : \italique{Paulys Realencyclopädie der classischen Altertumswissenschaft}, IV 1 (1931), coll. 1011-1012.}:\citComm{Συβήνη: ἡ δερματίνη αὐλοθήκη\nc{Dans certains manuscrits du \italique{Zonarae Lexicon} (p. 1686), on trouve βελοθήκη en place d'αὐλοθήκη.}. ἢ ἡ φαρέτρα. Θεσμοφοριαζούσαις Ἀριστοφάνης· ἀλλ' οὐκ ἔχω οὐδέν. ἀλλὰ συβήνην λάβε.}Cette unanimité ne doit pas nous tromper : elle prouve seulement l'ancienneté de la première rédaction qui aura diversement infusé dans de nombreux lexiques. Je dirais même plus : ces gloses concordent si bien qu'il est certain qu'elles remontent toutes à une même source -- certes ancienne -- mais unique, et donc faillible.\\
Que le terme συβήνη -- qui n'est d'ailleurs pas celui transmis par notre manuscrit comportant συ\textbf{μ}βήνην --, que συβήνη ait pu désigner un \italique{étui à flûte}, je n'en disconviens pas. À la rigueur, que par dérivation le mot ait aussi pu être appliqué à un \italique{carquois}, pourquoi pas\nc{Encore qu'il me semble plus vraisemblable que ce sens ait été inféré par notre scholiaste, puis scrupuleusement intégré à tous les lexiques.} ? Mais il serait absurde de lui donner ce sens dans notre passage : croit-on que l'Archer, tout Scythe qu'il fût, pensait un seul instant qu'une vieille maquerelle pût accepter son carquois dont elle n'aurait que faire au lieu d'une drachme sonnante et trébuchante ? ce serait tout bonnement ridicule\nc{À tous les vulgarisateurs d'Aristophane, tournebouleurs et avilisseurs de ses onze chefs-d'œuvre, rappelons la maxime de bon goût qu'Aristophane s'imposait :\citComm{ὡς ἐγὼ φοβοῦμαι περὶ τῶν μελλόντων ῥηθήσεσθαι, οὔ τι μὴ γελοῖα εἴπω — τοῦτο μὲν γὰρ ἂν κέρδος εἴη καὶ τῆς ἡμετέρας μούσης ἐπιχώριον — ἀλλὰ μὴ καταγέλαστα.\\
Ce n'est pas de provoquer le rire que je redoute -- cela serait tout gain, et conforme à ma muse -- mais de dire des choses ridicules.\\
(Trad. P. Boutang.)}(Tradition reçue en droite ligne d'Aristodème, témoin des événements, qui en fit le récit à Apollodore, qui lui-même le transmit à Glaucon, frère de Platon, qui -- on peut l'imaginer ainsi -- confia bientôt, après peut-être encore quelques méandres, à son frère Platon le précieux compte-rendu, lequel Platon finit par coucher par écrit cette fameuse nuit de \italique{Banquet}, 189b, quelques quarante années après les faits.)} ! Le Scythe, Euripide, Aristophane, le savent bien ; en fait n'importe quel homme, pourvu qu'il ne fût pas versé dans l'exégèse d'Aristophane : \citComm{Chez la belle Suzon, pas d'argent, pas de cuisse...\nc{Le proverbe d'origine détourné par Brassens (\italique{Grand-père}) étant «~point d'argent, point de Suisse~». Quel sera l'homme d'assez bonne volonté pour nous offrir le colossal équivalent brassinien du maître-œuvre de Taillardat, quel est le preux qui prendra à bras-le-corps les quelques deux cents chansons de Brassens que ce dernier a bien voulu nous conserver ? Espérons qu'il ne soit pas besoin d'attendre encore vingt-cinq siècles comme pour Aristophane. (C'est quand le vers est troussé qu'il faut l'analyser.) Je tiens à la disposition de cette bonne âme éventuelle un matériau qui poserait de solides fondations à ces \italique{Images de Brassens} en fomentation.}}En réalité, le Scythe demande à Euripide de garder son prisonnier quelques instants, le temps qu'il aille chercher l'argent nécessaire.
Encore une fois, c'est Hésychius (σ 2137) qui détenait la solution :\citComm{συμβίνης· καπροβόλον. ἐμβόλιον.}Le mot est apparenté à σιβύνη, une \italique{lance} (Hésychius nous précise pour la chasse au \italique{sanglier}, κάπρος), dont il est une légère déformation qui prépare le jeu de mots scabreux du vers 1215.
Le texte est aussi clair qu'un ciel d'azur au milieu de l'été : le Scythe confie sa lance à Euripide pour la garde de Mnésiloque.
Pour le texte, voir le commentaire critique.\\
Le Scythe part chercher l'argent, accompagné de la jeune beauté, qui se montre si caressante selon les instructions d'Euripide (vers 1172 ; l'attitude de la danseuse est la version muette de celle de la Myrrhine dans \italique{Lysistrata}) qu'il en revient tout émoustillé et en proie à une promesse de félicité qui point\nc{Savez-vous ce qu'on peut lire dans le commentaire le plus universitairement reçu, le plus académiquement académisant -- celui d'Austin \& Olson, puisqu'il faut le nommer ? Je vous le donne en mille : le Scythe a consommé la belle Élaphion... il revient donc, précisent leurs auteurs, avec l'\italique{olisbos} abaissé (!!!) Quand on sait -- ou qu'on ne sait pas, comme Austin \& Olson -- combien Aristophane adorait ses chères actrices, qu'il ne manquait jamais de rendre grâce à leur beauté, toujours plein d'attention et de prévenance vis-à-vis de la voracité de la lie athénienne... un archer Scythe... la lie de la lie... suggérer une telle monstruosité ?! «~Paraît qu'elle s'vend même à des flics. / Quelle décadence ! / Y'a plus d'moralité publique / Dans notre France...~» En France peut-être, mais au théâtre d'Aristophane, certainement pas ! (Lisez la note de Willems, t. III pp. 394-395, in « Le nu dans la Comédie ancienne des Grecs », article reproduit en annexe ; purgatif si nécessaire quand on est souillé par de telles salauderies !) A\&O ne nous indiquent pas si le public pouvait se délecter en sus de bruitages émanant des coulisses, quand ça y fornique... Mais pour qui prennent-ils Aristophane ?.. voici donc la littérature qu'on inflige à nos pauvres étudiants... je tiens le cas pour pendable !\citComm{et si, par un malheur, j'en avois fait autant,\\
je m'irois, de regret, pendre tout à l'instant.}}. Il déchante très vite : Euripide et Mnésiloque ont pris la poudre d'escampette. Élaphion décampe à ce moment sous les imprécations du Scythe (ἀπότρεκ᾽ ὠς τάκιστα σύ).
Le Scythe s'arme donc de sa lance, laissée sur place, et s'apprête à descendre (καταβαίνειν) à leur poursuite :\citComm{ὀρτῶς δὲ συμβήνη ᾽στί· καταβηνῆσι γάρ.}Je n'explique pas le trait d'esprit -- d'esprit \italique{Scythe} ; si vous ne le comprenez pas, laissez-là Aristophane et allez faire des mathématiques !}
\comm{1201}{Ἀρταμουξία.}{L'archer scythifie le nom d'Ἀρτεμισία : il l'entend comme un composé de ἀρτα-, multiplie les `α' et ajoute un `ξ'\nc{Voir le commentaire vv. 91-2 d'Olson aux \italique{Acharniens} ainsi que l'article «~The Archer Scene in Aristophanes' \italique{Thesmophoriazusae}~» d'Édith M. Hall, in \italique{Philologus} 133 (1989), p. 40. On peut citer Ἀρταβάτας, Ξέρξης, Ἀρταξέρξης, Ἀρταφρένης, Ἀρτάβαζος \&c.}. Cf. le parler de Ψευδ\textbf{αρτά}βας dans les \italique{Acharniens} (v. 100) :\citComm{ἰαρταμὰν ἐξάρξαν ἀπισσόνα σάτρα.}}
